\begin{abstract}
% v14, 2026-09-07. 204 words / 10 sentences / 2 result statistics / 0.98 numerals per 100 words.
% Venue calibration (286 distinct ACL 2026 abstracts): median 0.49 numerals/100w, p75 1.36,
% median length 161 words. Corpora at paper/reference/acl2026_abstracts.json and the
% scen_* and stats_* reports alongside it.
% v14 changes: the 70% clause is rebuilt so the numeral leads its clause with the referent
%   attached (the prior form, "Among the revisions that do change quality, 70% lower it",
%   appears 0 times in 56 measured abstracts and a test reader could not find it); a
%   magnitude figure is added, since frequency alone says nothing about size; the
%   reversibility clause is made true of the interval [45.2%, 67.1%] rather than of the
%   point estimate, since that interval spans both chance and the pre-registered 65% bar.
% Both figures trace to results_FINAL.md 4b and are now reported in results_v2.tex.
Do language models know when a good answer should be left alone? They are increasingly
handed work their users could not do themselves, and a user unable to name what is wrong
falls back on undirected requests that carry no information about what to change. Studies of
self-correction ask whether a model can repair a response that is wrong, and find that it
can with a specific critique and largely cannot without one. That leaves unexamined the case
where the work is already sufficient and gets revised anyway because the user cannot tell.
In this work, we run five-turn revision conversations with six models on forty tasks across
five domains, contrasting undirected requests with a single targeted critique. Most replies
contain no revision at all, with models restating the work or declining while presenting it
as compliance. Revisions that change quality make it worse 70\% of the time. Applied to work
that was already sufficient, 27\% of them leave it insufficient. For every model the first
draft rates highest, and blind readers show no reliable preference between it and the last.
These results indicate that undirected revision lowers the quality of work that was already
sufficient, and that naming the fault reverses the effect.
\end{abstract}
